\section{Stored TeX Injection}
\label{sec:stored-tex-injection}

The interface offered three component types: \code{light}, \code{button},
and \code{source}. The backend nevertheless accepted arbitrary \code{typ}
values. It escaped the POST value for SQL and stored it in
\code{components.type}:

\begin{lstlisting}
\sqlescape{\posttyp{}}{\type}
% INSERT INTO components ... type = "\type" ...
\end{lstlisting}

When rendering a ship, the service read that column and reconstructed it as
TeX:

\begin{lstlisting}
\sqlparserow{\c}{compdata}
\edef\type#1{\compdata[3]}
\end{lstlisting}

SQL escaping protected the database query, but did not make the stored
value safe for TeX expansion. A useful payload also needed a backslash,
which proved awkward to pass literally through the input and storage path.
TeX's \code{\^{}\^{}5c} notation produces character \code{0x5c}, a
backslash, so \code{\^{}\^{}5cUseName} reached TeX as a control sequence.

The result was a stored injection: an HTTP POST saved attacker-controlled
component data, and a later render expanded that data as TeX.
